%==============================================================================
% Appendix: Proof of Theorem 5 — EntroDiff Reverse Sampler as JKO Gradient Flow
% Auto-converted from Docs/proof/Theorem 5 draft.md
%==============================================================================
\begin{theorem}[JKO correspondence]
\label{appx-thm:jko}
Let $\ptau$ denote the density of the reverse sampler at diffusion time
$\tau \in [0,T]$, and let $\rhotrue$ be the target distribution induced by the PDE
solution in data space.
Define the free energy functional
\begin{equation}
\mathcal{F}(\rho) = \Ent(\rho \mid \rhotrue) + \lambda\,\mathcal{R}(\rho),
\qquad
\Ent(\rho \mid \rhotrue) = \int_{\R^{N}} \rho(x)\,\log\frac{\rho(x)}{\rhotrue(x)}\,dx,
\end{equation}
where $\mathcal{R}$ is a density functional encoding the PDE constraint and $\lambda \ge 0$.

Assume the following hypotheses:
\begin{enumerate}
\item \textbf{(Exact score limit.)}
  After sufficient training, $\sth(x,\tau) = \nabla_x \log \ptau(x)$ pointwise.

\item \textbf{(Viscosity-matched schedule.)}
  $a(\tau) = g(\tau)^{2}/2 > 0$, aligned with the physical viscosity.

\item \textbf{(PDE constraint as a differentiable functional.)}
  $\mathcal{R}$ is Fr\'{e}chet-differentiable on $\mathcal{P}_{2}(\R^{N})$, with
  first variation $\delta \mathcal{R} / \delta \rho$ well-defined;
  $\nabla_x (\delta \mathcal{R} / \delta \rho)$ is locally bounded.

\item \textbf{(Rigid $\Lburg$ consistency.)}
  The Burgers consistency penalty~$\Lburg$ is sufficiently strong that, in the
  infinite-capacity noiseless limit, the learned score field satisfies the viscous
  Burgers structure
  \begin{equation}
  \dtau \score_{\tau} + 2\,\score_{\tau}\cdot\nabla\,\score_{\tau}
    = \Lap \score_{\tau}.
  \end{equation}
  Moreover, $\ptau \in \mathcal{P}_{2}(\R^{N})$ is smooth in $(x,\tau)$ and strictly
  positive ($\ptau > 0$).

\item \textbf{(PDE guidance as a gradient drift.)}
  In the continuous-time limit, the combined PDE constraint losses $\Lent$
  (Kruzhkov entropy penalty) and $\Lburg$ (Burgers consistency) induce an additional
  guidance velocity $g_{\tau}(x)$ in the reverse ODE that converges to the spatial
  gradient of a functional.
  Specifically, the total effective velocity field of the EntroDiff reverse sampler is
  \begin{equation}
  \label{eq:total-vel-stmt}
  v_{\tau}(x)
    = -a(\tau)\Bigl[
        \nabla_x\log\ptau(x)
        + \nabla_x V(x)
        + \lambda\,\nabla_x\frac{\delta\mathcal{R}}{\delta\rho}(\ptau)(x)
      \Bigr],
  \end{equation}
  where $V$ is the potential of the target distribution ($\rhotrue \propto e^{-V}$).
  The guidance contributions $\nabla V$ and $\lambda\nabla(\delta\mathcal{R}/\delta\rho)$
  arise from the gradient of the PDE losses with respect to the sample trajectory,
  which by design converge to these functional-gradient limits when the losses are driven
  to zero in training.
\end{enumerate}

Then the density evolution of the reverse sampler satisfies
\begin{equation}
\label{eq:density-flow}
\dtau \ptau = \nabla \cdot \Bigl( a(\tau)\, \ptau\, \nabla\frac{\delta\mathcal{F}}{\delta\rho}(\ptau) \Bigr).
\end{equation}
Equivalently, under the time reparameterization
$\theta(\tau) = \int_{0}^{\tau} a(r)\,dr$, we have
\begin{equation}
\label{eq:W2-flow}
\partial_{\theta} p_{\theta}
  = \nabla \cdot \Bigl( p_{\theta}\, \nabla\frac{\delta\mathcal{F}}{\delta p}(p_{\theta}) \Bigr),
\end{equation}
which is precisely the $\Wass{2}$-gradient flow of~$\mathcal{F}$.
Consequently, by the JKO theory, a single time-discrete step takes the variational form
\begin{equation}
\label{eq:jko-step}
p^{n+1}
  = \arg\min_{p \in \mathcal{P}_{2}(\R^{N})}
    \Bigl\{ \mathcal{F}(p) + \frac{1}{2\Delta\theta}\,\Wass{2}^{2}(p, p^{n}) \Bigr\}.
\end{equation}
\end{theorem}

\begin{proof}

%==========================================================================
% Step 1: Reverse Sampler as a Continuity Equation
%==========================================================================
\medskip
\noindent\textbf{Step 1: Reverse sampler as a continuity equation.}

Let $X_{\tau}$ denote the particle trajectory of the reverse sampler.
If its continuous-time limit satisfies the ODE
\begin{equation}
\frac{d X_{\tau}}{d\tau} = v_{\tau}(X_{\tau}),
\end{equation}
then by the standard pushforward formula, the particle distribution
$\ptau = \Law(X_{\tau})$ satisfies the continuity equation
\begin{equation}
\label{eq:continuity}
\dtau \ptau + \nabla \cdot (\ptau\,v_{\tau}) = 0.
\end{equation}
Thus, identifying the velocity field $v_{\tau}$ is the central task.

For any test function $\varphi \in C_{c}^{\infty}(\R^{N})$, the identity follows from the
chain rule:
\begin{align}
\frac{d}{d\tau} \int_{\R^{N}} \varphi(x)\,\ptau(x)\,dx
  &= \frac{d}{d\tau}\,\E[\varphi(X_{\tau})] \notag\\
  &= \E[\nabla\varphi(X_{\tau}) \cdot v_{\tau}(X_{\tau})] \notag\\
  &= \int \nabla\varphi(x) \cdot v_{\tau}(x)\,\ptau(x)\,dx.
\end{align}
Integration by parts of the right-hand side yields
\begin{equation}
\frac{d}{d\tau} \int \varphi\,\ptau\,dx
  = - \int \varphi(x)\,\nabla\cdot(\ptau\,v_{\tau})(x)\,dx,
\end{equation}
and since $\varphi$ is arbitrary, the continuity equation~\eqref{eq:continuity} holds in
the sense of distributions (and pointwise under the smoothness of $\ptau$ and $v_{\tau}$).

%==========================================================================
% Step 2: Score Drift and PDE Guidance — Total Effective Velocity
%==========================================================================
\medskip
\noindent\textbf{Step 2: Score drift and PDE guidance --- the total effective velocity.}

We identify $v_{\tau}$ as the sum of two contributions.

\smallskip
\noindent\textit{(a)  Score drift (from denoising score matching).}
Under Hypothesis~1 (exact score) and Hypothesis~4 (rigid $\Lburg$), the network score
converges to the true score:
\begin{equation}
\sth(x,\tau) = \nabla_x \log \ptau(x).
\end{equation}
The probability flow ODE induced by score matching alone has the velocity
\begin{equation}
\label{eq:score-vel}
v_{\tau}^{\mathrm{score}}(x) = -a(\tau)\,\sth(x,\tau) = -a(\tau)\,\nabla_x\log\ptau(x).
\end{equation}

\smallskip
\noindent\textit{(b)  PDE guidance drift (from constraint losses).}
The EntroDiff loss family includes PDE-constraint terms beyond denoising score matching:
$\Lent$ (Kruzhkov entropy penalty) and $\Lburg$ (Burgers consistency).
These losses are functionals of the generated trajectory that penalize deviations from
physical admissibility.
In the continuous-time and infinite-data limit, their gradients with respect to the sample
position induce an additional drift $g_{\tau}(x)$.

By Hypothesis~5, as the constraint losses are driven to zero, this guidance drift converges
to the gradient of the potential that defines the target distribution plus the variational
gradient of the PDE-constraint functional:
\begin{equation}
\label{eq:guidance-vel}
g_{\tau}(x) = - \Bigl( \nabla_x V(x) + \lambda\,\nabla_x\frac{\delta\mathcal{R}}{\delta\rho}(\ptau)(x) \Bigr).
\end{equation}
Here $\nabla V = -\nabla\log\rhotrue$ is the drift toward the target $\rhotrue \propto e^{-V}$,
and $\lambda\nabla(\delta\mathcal{R}/\delta\rho)$ is the PDE-constraint drift
(see Step~4 for the construction of~$\mathcal{R}$).

\smallskip
\noindent\textit{(c)  Total effective velocity.}
Summing~\eqref{eq:score-vel} and~\eqref{eq:guidance-vel}, and absorbing the common
factor $-a(\tau)$,
\begin{equation}
\label{eq:total-vel}
v_{\tau}(x)
  = -a(\tau)\Bigl[
      \nabla_x\log\ptau(x)
      + \nabla_x V(x)
      + \lambda\,\nabla_x\frac{\delta\mathcal{R}}{\delta\rho}(\ptau)(x)
    \Bigr].
\end{equation}
Substituting into the continuity equation gives
\begin{equation}
\label{eq:cont-with-vel}
\dtau \ptau
  = \nabla \cdot \Bigl(
      a(\tau)\,\ptau\,
      \bigl[ \nabla\log\ptau + \nabla V + \lambda\nabla\tfrac{\delta\mathcal{R}}{\delta\rho} \bigr]
    \Bigr).
\end{equation}
We now show that the bracketed quantity in~\eqref{eq:total-vel} is exactly
$\nabla(\delta\mathcal{F}/\delta\rho)$.

%==========================================================================
% Step 3: First Variation of the Entropy Term
%==========================================================================
\medskip
\noindent\textbf{Step 3: First variation of the entropy term.}

Consider the relative entropy
\begin{equation}
\Ent(p \mid \rhotrue) = \int p\,\log\frac{p}{\rhotrue}\,dx.
\end{equation}
Take a perturbation $p_{\varepsilon} = p + \varepsilon\eta$ with $\int\eta\,dx = 0$.
A first-order expansion gives
\begin{equation}
\frac{d}{d\varepsilon}\Big|_{\varepsilon=0} \Ent(p_{\varepsilon} \mid \rhotrue)
  = \int \eta(x)\Bigl( \log\frac{p(x)}{\rhotrue(x)} + 1 \Bigr)\,dx.
\end{equation}
Hence the $L^{2}$ first variation is
\begin{equation}
\frac{\delta}{\delta p}\,\Ent(p \mid \rhotrue) = \log\frac{p}{\rhotrue} + 1
\end{equation}
(the additive constant does not affect the gradient flow).
Taking the spatial gradient,
\begin{equation}
\label{eq:ent-grad}
\nabla\frac{\delta}{\delta p}\,\Ent(p \mid \rhotrue)
  = \nabla\log p - \nabla\log\rhotrue.
\end{equation}
Since $\rhotrue \propto e^{-V}$, we have $-\nabla\log\rhotrue = \nabla V$, and therefore
\begin{equation}
\label{eq:ent-var}
\nabla\frac{\delta}{\delta p}\,\Ent(p \mid \rhotrue) = \nabla\log p + \nabla V.
\end{equation}

%==========================================================================
% Step 4: PDE Constraint as a Differentiable Functional
%==========================================================================
\medskip
\noindent\textbf{Step 4: PDE constraint as a differentiable functional.}

The PDE constraint (beyond $\Lburg$) must be expressible as a density
functional~$\mathcal{R}(\rho)$.
Under Hypothesis~3, $\mathcal{R}$ is Fr\'{e}chet-differentiable on $\mathcal{P}_{2}$, so
its first variation $\frac{\delta\mathcal{R}}{\delta\rho}$ exists and satisfies
\begin{equation}
\frac{d}{d\varepsilon}\Big|_{\varepsilon=0} \mathcal{R}(\rho_{\varepsilon})
  = \int \frac{\delta\mathcal{R}}{\delta\rho}(\rho)\,\eta\,dx
\end{equation}
for any zero-mass perturbation $\eta$.

In the Wasserstein setting, perturbations are generated by a transport
field~$\xi$ via the pushforward:
\begin{equation}
T_{\varepsilon}(x) = x + \varepsilon\,\xi(x), \qquad
\rho_{\varepsilon} = (T_{\varepsilon})_{\#}\,\rho.
\end{equation}
The corresponding variational formula is
\begin{equation}
\frac{d}{d\varepsilon}\Big|_{\varepsilon=0}
  \mathcal{R}\bigl((\mathrm{Id} + \varepsilon\xi)_{\#}\rho\bigr)
  = \int \nabla\frac{\delta\mathcal{R}}{\delta\rho}(\rho) \cdot \xi\;\rho\,dx.
\end{equation}
Thus, the $\mathcal{R}$ gradient flow term contributes a drift
$-\nabla(\delta\mathcal{R}/\delta\rho)$ in the velocity field (see~\eqref{eq:guidance-vel}).

%==========================================================================
% Step 5: Identification with the Free Energy Gradient
%==========================================================================
\medskip
\noindent\textbf{Step 5: Identification with the free energy gradient.}

Combining the entropy and PDE-constraint terms, the free energy is
\begin{equation}
\mathcal{F}(\rho) = \Ent(\rho \mid \rhotrue) + \lambda\,\mathcal{R}(\rho).
\end{equation}
From~\eqref{eq:ent-var} and Step~4, its spatial variational gradient is
\begin{equation}
\label{eq:F-grad}
\nabla\frac{\delta\mathcal{F}}{\delta\rho}
  = \nabla\frac{\delta}{\delta\rho}\,\Ent(\rho \mid \rhotrue)
    + \lambda\,\nabla\frac{\delta\mathcal{R}}{\delta\rho}
  = \nabla\log\rho + \nabla V + \lambda\,\nabla\frac{\delta\mathcal{R}}{\delta\rho}.
\end{equation}

Comparing~\eqref{eq:total-vel} with~\eqref{eq:F-grad}, the total effective velocity is exactly
\begin{equation}
\label{eq:vel-F}
v_{\tau}(x) = -a(\tau)\,\nabla_x\frac{\delta\mathcal{F}}{\delta\rho}(\ptau)(x).
\end{equation}

Substituting into the continuity equation (Step~1),
\begin{equation}
\dtau \ptau + \nabla \cdot \Bigl( \ptau\Bigl[ -a(\tau)\,\nabla\frac{\delta\mathcal{F}}{\delta\rho}(\ptau) \Bigr] \Bigr) = 0,
\end{equation}
i.e.,
\begin{equation}
\label{eq:main-pde}
\dtau \ptau = \nabla \cdot \Bigl( a(\tau)\,\ptau\,\nabla\frac{\delta\mathcal{F}}{\delta\rho}(\ptau) \Bigr).
\end{equation}

To eliminate the time-weight $a(\tau)$, define the reparameterized time
\begin{equation}
\theta(\tau) = \int_{0}^{\tau} a(r)\,dr.
\end{equation}
Since $a(\tau) > 0$, $\theta$ is strictly increasing and locally invertible.
By the chain rule,
\begin{equation}
\dtau \ptau = \frac{d\theta}{d\tau}\,\partial_{\theta}p_{\theta}
            = a(\tau)\,\partial_{\theta}p_{\theta}.
\end{equation}
Cancelling $a(\tau)$ yields the canonical Wasserstein-2 gradient flow form
\begin{equation}
\label{eq:W2-final}
\partial_{\theta} p_{\theta}
  = \nabla \cdot \Bigl( p_{\theta}\,\nabla\frac{\delta\mathcal{F}}{\delta p}(p_{\theta}) \Bigr).
\end{equation}

This completes the bridge from the EntroDiff reverse sampler to an exact Wasserstein
gradient flow. $\square$

%==========================================================================
% Step 6: JKO Discretization from the Wasserstein Gradient Flow
%==========================================================================
\medskip
\noindent\textbf{Step 6: JKO discretization from the Wasserstein gradient flow.}

We now show that the implicit Euler discretization of the above gradient flow is equivalent to
the JKO variational scheme.

\smallskip
\noindent\textit{6.1  Implicit Euler discretization.}
Consider a uniform time grid with step $\Delta\theta$ (in reparameterized time).
Given $p^{n} \approx p_{\theta_{n}}$, the implicit Euler discretization of the gradient flow is
\begin{equation}
\label{eq:implicit-euler}
\frac{p^{n+1} - p^{n}}{\Delta\theta}
  = \nabla \cdot \Bigl( p^{n+1}\,\nabla\frac{\delta\mathcal{F}}{\delta\rho}(p^{n+1}) \Bigr).
\end{equation}

\smallskip
\noindent\textit{6.2  Proximal step analogy.}
In Euclidean space, the implicit Euler step for $\dot x = -\nabla F(x)$ is equivalent to the
proximal minimization
\begin{equation}
x^{n+1} = \arg\min_{x} \Bigl\{ F(x) + \frac{1}{2\Delta t}\,\|x - x^{n}\|^{2} \Bigr\}.
\end{equation}
The Wasserstein analogue replaces the Euclidean distance with $\Wass{2}$ and the Euclidean
gradient with the continuity equation structure.
The claim is that
\begin{equation}
\label{eq:jko-claim}
p^{n+1}
  = \arg\min_{p \in \mathcal{P}_{2}(\R^{N})}
    \Bigl\{ \mathcal{F}(p) + \frac{1}{2\Delta\theta}\,\Wass{2}^{2}(p, p^{n}) \Bigr\}
\end{equation}
is the JKO step corresponding to the implicit Euler discretization above.

\smallskip
\noindent\textit{6.3  Variational derivation.}
Define the augmented functional
\begin{equation}
\mathcal{J}(p) = \mathcal{F}(p) + \frac{1}{2\Delta\theta}\,\Wass{2}^{2}(p, p^{n}),
\end{equation}
and let $p^{n+1}$ be a minimizer.
In Wasserstein space, admissible variations are generated by transport maps rather than
additive perturbations.
Define the perturbation
\begin{equation}
T_{\varepsilon}(x) = x + \varepsilon\,\xi(x), \qquad
p_{\varepsilon} = (T_{\varepsilon})_{\#}\,p^{n+1},
\end{equation}
where $\xi$ is an arbitrary smooth, compactly supported vector field.
Optimality requires
\begin{equation}
\frac{d}{d\varepsilon}\Big|_{\varepsilon=0} \mathcal{J}(p_{\varepsilon}) = 0.
\end{equation}

\textbf{Variation of $\mathcal{F}$.}
By the Wasserstein variational formula,
\begin{equation}
\frac{d}{d\varepsilon}\Big|_{\varepsilon=0} \mathcal{F}(p_{\varepsilon})
  = \int \nabla\frac{\delta\mathcal{F}}{\delta\rho}(p^{n+1}) \cdot \xi\;
    p^{n+1}\,dx.
\end{equation}

\textbf{Variation of the $\Wass{2}^{2}$ term.}
Let $\pi$ be the optimal transport plan from $p^{n+1}$ to $p^{n}$, with Kantorovich
potential~$\phi$.
A classical result of optimal transport theory gives
\begin{equation}
\frac{\delta}{\delta\rho}\Bigl( \frac{1}{2}\,\Wass{2}^{2}(\rho, p^{n}) \Bigr) = \phi,
\end{equation}
and therefore
\begin{equation}
\frac{d}{d\varepsilon}\Big|_{\varepsilon=0} \frac{1}{2}\,\Wass{2}^{2}(p_{\varepsilon}, p^{n})
  = \int \nabla\phi \cdot \xi\;
    p^{n+1}\,dx.
\end{equation}

\textbf{First-order condition.}
Combining both terms,
\begin{equation}
\int \Bigl[ \nabla\frac{\delta\mathcal{F}}{\delta\rho}(p^{n+1})
          + \frac{1}{\Delta\theta}\,\nabla\phi \Bigr] \cdot \xi\;
      p^{n+1}\,dx = 0.
\end{equation}
Since $\xi$ is arbitrary, the integrand must vanish $p^{n+1}$-almost everywhere, yielding
\begin{equation}
\label{eq:nabla-phi}
\nabla\phi = -\Delta\theta\,\nabla\frac{\delta\mathcal{F}}{\delta\rho}(p^{n+1}).
\end{equation}

\textbf{Connection to the optimal transport map.}
The Brenier theorem gives the optimal transport map $T$ from $p^{n+1}$ to $p^{n}$ as
$T(x) = x - \nabla\phi(x)$.
Substituting the expression for $\nabla\phi$,
\begin{equation}
T(x) = x + \Delta\theta\,\nabla\frac{\delta\mathcal{F}}{\delta\rho}(p^{n+1}).
\end{equation}

\textbf{Deriving the discrete continuity equation.}
By the pushforward relation $p^{n} = T_{\#}p^{n+1}$, we expand for small $\Delta\theta$
(using the Jacobian expansion of the pushforward):
\begin{equation}
p^{n}(x)
  = p^{n+1}(x)
    - \Delta\theta\,\nabla\cdot\Bigl( p^{n+1}\,\nabla\frac{\delta\mathcal{F}}{\delta\rho}(p^{n+1}) \Bigr)
    + o(\Delta\theta).
\end{equation}
Rearranging,
\begin{equation}
\frac{p^{n+1} - p^{n}}{\Delta\theta}
  = \nabla\cdot\Bigl( p^{n+1}\,\nabla\frac{\delta\mathcal{F}}{\delta\rho}(p^{n+1}) \Bigr) + o(1).
\end{equation}
Taking $\Delta\theta \to 0$, we recover exactly the implicit Euler discretization~\eqref{eq:implicit-euler}
of the Wasserstein gradient flow.

\smallskip
\noindent\textit{6.4  Conclusion of the JKO derivation.}
Thus, a minimizer $p^{n+1}$ of the JKO functional satisfies the implicit Euler discretization
of the Wasserstein gradient flow.
The reverse sampler, whose density evolution was established in Steps~1--5 as the gradient
flow of~$\mathcal{F}$, therefore admits the JKO variational characterization in the
discrete-time limit $\Delta\tau \to 0$.
This completes the proof of Theorem~5.
\end{proof}

%==========================================================================
% Compact Fokker--Planck Expansion
%==========================================================================
\medskip
\noindent\textbf{Compact Fokker--Planck expansion.}

For completeness, expanding the gradient flow structurally: using
$\nabla\log p_{\theta} = \nabla p_{\theta}/p_{\theta}$ and the explicit form of the
entropy variation from Step~3, the gradient flow PDE becomes
\begin{equation}
\begin{aligned}
\partial_{\theta} p_{\theta}
  &= \Lap p_{\theta}
     + \nabla\cdot(p_{\theta}\nabla V)
     + \lambda\,\nabla\cdot\Bigl( p_{\theta}\,\nabla\frac{\delta\mathcal{R}}{\delta\rho} \Bigr),
\end{aligned}
\end{equation}
exhibiting the standard decomposition into a diffusion term ($\Lap p_{\theta}$),
an external potential drift ($\nabla\cdot(p_{\theta}\nabla V)$),
and a PDE-constraint transport term.

%==========================================================================
% Remark: Relation to the Obukhov--Fokker--Planck Correspondence
%==========================================================================
\medskip
\noindent\textbf{Remark (Relation to the Obukhov--Fokker--Planck correspondence).}

The identification~\eqref{eq:vel-F} can be understood through the lens of
entropy-dissipation balances.
The pure reverse probability-flow ODE ($v = -a\nabla\log p$) is the Wasserstein gradient flow
of the negative entropy $H(p) = \int p\log p$.
Adding the PDE guidance terms $\nabla V + \lambda\nabla(\delta\mathcal{R}/\delta\rho)$ shifts
the stationary point from the standard Gaussian to the target~$\rhotrue$ while incorporating
the physical constraint~$\mathcal{R}$.
The resulting free energy $\mathcal{F} = \Ent(\cdot\mid\rhotrue) + \lambda\mathcal{R}$ is
the natural Lyapunov functional for the guided dynamics, and the JKO scheme is its canonical
time discretization~\citep{jko1998}.
